
\subsection{Proofs for Section~\ref{sec:exactness}}\label{app:exactness-proofs}

In Propositions~\ref{prop:one-probe}, \ref{prop:zero-probe-fermat},
\ref{prop:separation}(c) and~\ref{prop:one-state} the register has one
value, so $L_m=1$ satisfies every register condition. Every table entry
is positive ($\theta_m\ge\delta>0$, $Q_m>0$), and feasibility at a profile
means the probability floor, normalization, the budget and the held
tests. Without probes every base kernel is
uniform, so $P_m((g,r),(h,r'))=\theta_m\,L_m(g,h;r,r')/|X|$, and by
Lemma~\ref{lem:compile} whether a parameter vector passes a held test
depends only on~$\theta$.

% ---------------------------------------------------------------------------
\subsubsection{Real versus rational reachability}\label{app:real-rational}

\begin{proof}[Proof of Proposition~\ref{prop:one-probe}]
By Lemma~\ref{lem:compile}(b), the word-against-pair test gives
$\theta_{\mathsf u}=\theta_{\mathsf x_4}+\theta_{\mathsf y_4}\ge2/9$. The
other seven labels have probability at least $1/9$, so normalization forces
$\theta_{\mathsf u}=2/9$ and all others to be $1/9$. By
Lemma~\ref{lem:compile}(a), the other probe equations give
$\alpha_{\mathsf x_2}=4x^2$, $\alpha_{\mathsf x_4}=4\alpha_{\mathsf
x_2}^2=64x^4$, $\alpha_{\mathsf y_2}=4y^2$ and $\alpha_{\mathsf
y_4}=64y^4$. The probe equation of the word-against-pair test reads
$(2/9)(1/8)=(64/9)(x^4+y^4)$, or $x^4+y^4=4^{-4}$. Conversely, each
point of this quartic gives eigenvalues in $[-1/4,1/4]$, within the
budget, passing every held test. The
points $((1-\sigma)^{1/4}/4,\sigma^{1/4}/4)$, $0\le\sigma\le1$, give the
real path. By Fermat's theorem for exponent four, applied to
$(4x)^4+(4y)^4=1$, the feasible configurations are the four with $(x,y)$
in $\{(\pm1/4,0),(0,\pm1/4)\}$. With the probe~$\chi$, the diagonal entry
is $P_{\mathsf x}((g,r),(g,r))=\theta_{\mathsf x}(1+x)/2=(1+x)/18$. Sequences
stay at the profile (Remark~\ref{rem:inputs}(ii)), and a step leaving $x=1/4$ for
$x\in\{0,-1/4\}$ changes this entry by at least $1/72$.
\end{proof}

\begin{proof}[Proof of Proposition~\ref{prop:zero-probe-fermat}]
By Lemma~\ref{lem:compile}, the held tests read
\[
  \theta_{\mathsf u}=\theta_{\mathsf x_4}+\theta_{\mathsf y_4},\qquad
  \theta_{\mathsf x_4}\theta_{\mathsf h}^3=\theta_{\mathsf x}^4,\qquad
  \theta_{\mathsf y_4}\theta_{\mathsf h}^3=\theta_{\mathsf y}^4,\qquad
  \theta_{\mathsf u}=17\,\theta_{\mathsf h}.
\]
Dividing by $\theta_{\mathsf h}>0$ and normalizing gives $x^4+y^4=17$ and
\[
  (\theta_{\mathsf u},\theta_{\mathsf h},\theta_{\mathsf x},\theta_{\mathsf
  y},\theta_{\mathsf x_4},\theta_{\mathsf y_4})=(17,1,x,y,x^4,y^4)/(35+x+y).
\]
Conversely, positive rational solutions give values passing every held
test, feasible exactly when each is at least $1/40$.
At $(1,2)$ and $(2,1)$ the normalizer is $38$ and the least numerator
is~$1$, so the least label probability is $1/38$, and $A$ and $B$ exist
and are unique. Along $1\le x\le2$ with $y=(17-x^4)^{1/4}$, also
$1\le y\le2$, so the normalizer is at most $39$, each numerator is at
least~$1$, every label probability exceeds $1/40$, and these real
parameter vectors form the required path.

Finally, Proposition~\ref{prop:zero-probes}(b) applies with
$w_0=\mathsf u$, $w_1=\mathsf x_4$ and $w_2=\mathsf y_4$. A vector $v$
lies in $\ker\Psi$ exactly when $v_{\mathsf u}=v_{\mathsf h}$,
$v_{\mathsf x_4}=4v_{\mathsf x}-3v_{\mathsf h}$ and $v_{\mathsf
y_4}=4v_{\mathsf y}-3v_{\mathsf h}$, so $(f\cdot v,g\cdot v)=(4(v_{\mathsf
x}-v_{\mathsf h}),4(v_{\mathsf y}-v_{\mathsf h}))$ and
$(f,g)(\ker_{\mathbb Z}\Psi)=4\mathbb Z^2$. Hence $d_0=d_1=d_\infty=4$,
with $3/4<1$, while $t=x^4/17$ is $1/17$ at $A$ and $16/17$ at~$B$. So $B$
is not finely reachable from~$A$.
\end{proof}

% ---------------------------------------------------------------------------
\subsubsection{The separation cutoff}\label{app:separation}

\begin{proof}[Proof of Proposition~\ref{prop:separation}]
Write $\Pi=n!\,D_0^{\,n}$, so $D=1+D_1n\Pi+D_1n^3\Pi^2D_0$. By the choice
of~$D_1$, $2^{-D_1}\le|c|\le2^{D_1}$ for every nonzero coefficient~$c$ and
$\kappa_\beta\le2^{D_1}$ for every budget.

\emph{Radius.} Let $x\in\FQ$ have support~$S$, let $G\subseteq S$ be the
coordinates lying in some budget, and put $\lambda_j=\log_2|x_j|$
for $j\in S$. Each budget meeting $S$ has $\kappa_\beta>0$, so
$\lambda_j\le D_1$ on~$G$. A test active on $S$ has both monomials nonzero
at~$x$, so its coefficient $c_\iota$ is nonzero and
$(\Delta_S\lambda)_\iota=\log_2|c_\iota|$ has absolute value at
most~$D_1$. Fix $i\in S$. An integral $w\in\RR^S$ with $\Delta_Sw=0$,
$w_j\le0$ on $G$ and $w_i\ge1$ would give rational points with coordinates
$x_j\varsigma^{-w_j}$ on $S$ and $0$ off~$S$. As rational
$\varsigma\downarrow0$, they satisfy every test (as in the proof of
Lemma~\ref{lem:restriction}(a)) and every budget but have $|x_i|\to\infty$,
against boundedness. %
Any real direction with $w_i>0$ could be scaled to satisfy $w_i\ge1$. The
resulting rational linear system would then have a rational solution, and
clearing denominators would give an integral direction. Hence no such real
$w$ with $w_i>0$ exists. By Farkas' lemma, $e_i$ is a linear combination of
the vectors $a-b$ of the tests active on~$S$ (restricted to~$S$), plus a
nonnegative combination of the $e_j$ with $j\in G$.
By Carath\'eodory's theorem, after splitting coefficients by sign, at
most $n$ linearly independent generators suffice, and Cramer's rule bounds
each coefficient by~$\Pi$. The inner product with $\lambda$ gives
$\lambda_i\le nD_1\Pi$, so $|x_i|\le2^{nD_1\Pi}$ on~$\FQ$.

\emph{Survivors.} Let $C=\core(S)\neq\emptyset$. Iterating
Lemma~\ref{lem:restriction}(a) gives $\hat x\in\FQ$ with support $C$ that
agrees with $x$ on~$C$. By Theorem~\ref{thm:certificates}(b), $C$ has a
Gordan certificate $\mu$ with at most $n$ nonzero multipliers, each of
absolute value at most~$\Pi$, so $v=\Delta_C^{\mathsf T}\mu$ is integral
with $1\le v_i\le n\Pi D_0$ on~$C$. Taking logarithms at $\hat x$ of the
tests with nonzero multipliers (coefficients $c_\iota$) gives
\[
  \sum_{i\in C}v_i\lambda_i=\sum_\iota\mu_\iota\log_2|c_\iota|\ge-nD_1\Pi .
\]
For each $i\in C$, bounding the other terms by the radius gives
$v_i\lambda_i\ge-nD_1\Pi-n\cdot n\Pi D_0\cdot nD_1\Pi=-(D-1)$; since
$v_i\ge1$, $|x_i|\ge2^{1-D}$ on $\core(\supp x)$.

\emph{(a)} Let $x,x'\in\FQ$ with $\|x-x'\|_\infty<2^{-D}$. Each
$i\in\core(\supp x)$ has $|x_i|\ge2^{1-D}>|x_i-x'_i|$, so $x'_i$ is
nonzero with the sign of~$x_i$, and symmetrically. Both cores lie in
$V=\supp x\cap\supp x'$, with the same signs. By
Lemma~\ref{lem:restriction}(b), $V$ peels to $V\cap\core(\supp
x)=\core(\supp x)$, which is irreducible (Lemma~\ref{lem:irreducible}), so
it is $\core(V)$ by Lemma~\ref{lem:restriction}(c); likewise for~$x'$. So
$\SC(x)=\SC(x')$; empty cores need no estimate.
For the size, $\Pi\le(n+1)^nD_0^n$ and $D\le3D_1(n+1)^3\Pi^2D_0$, with
$n,\log_2D_0\le\mathcal L$ and $D_1\le\mathcal L+1$, so $\log_2D$ is polynomial in~$\mathcal L$;
and $2^D$ has $D+1$ binary digits.

\emph{(b)} Steps never lower the profile, so along a sequence from $A$
each eigenvalue vector is a feasible rational point of the bounded system
$\Phi_{p_A}$ (Lemma~\ref{lem:compile}(a)). Profile steps keep this vector
fixed (Definition~\ref{def:profile-steps}), whereas numerical steps move
it by less than $2^{-D}$ (Lemma~\ref{lem:transfer}(a), as in the necessity
proof of Theorem~\ref{thm:fixed-profile} with $\varepsilon_\alpha=2^{-D}$).
By~(a), $\SC_{p_A}$ never changes.

\emph{(c)} The feasible rational points have $x_i=2^{-2^i}$ and
$z=\pm2^{-2^s}$, at distance $2\cdot2^{-2^s}$. On the full support, the
calibration $x_0=1/2$ forces $\eta_{x_0}=0$, then
$\eta_{x_{i+1}}=2\eta_{x_i}=0$ and $2\eta_z=2\eta_{x_s}=0$. Thus the full
support is irreducible and the signed cores differ at~$z$, while the
encoding has length polynomial in~$s$. The model uses a binary base with
one probe and one register value, self-reversed labels
$\mathsf h,\mathsf x_0,\dots,\mathsf x_s,\mathsf z$, $\delta=1/(s+3)$,
$\kappa=3/4$ and $H=2$. Its full-support level-$0$ profile holds
$\alpha_{\mathsf h}=1/2$, $\alpha_{\mathsf x_0}=1/4$ and the word tests
$\mathsf x_{i+1}\mathsf h\sim\mathsf x_i\mathsf x_i$ $(0\le i<s)$ and
$\mathsf{zz}\sim\mathsf x_s\mathsf x_s$, which form the only admissible
set. The probability floor forces every label probability to be $1/(s+3)$,
so the probe equations give $\alpha_{\mathsf x_{i+1}}/2=\alpha_{\mathsf
x_i}^2$, whence $\alpha_{\mathsf x_i}=x_i/2$ and $\alpha_{\mathsf z}=z/2$
for the two points. Both points are feasible. The corresponding weight equations force all
weights to be zero on the full support, so the signed cores differ
at~$\mathsf z$. Their tables differ only for~$\mathsf z$, by
$\theta_{\mathsf z}\cdot2\cdot2^{-2^s-1}/2=2^{-2^s-1}/(s+3)$.
\end{proof}

\emph{A binary-exponent variant.}
With binary exponents, the three-coordinate system $x_0=1/2,\
x_1=x_0^{2^b},\ z^2=x_1^2$, with the same coordinate budgets, has encoding
length $O(b)$ excluding the endpoints. Its two points have different
signed cores and gap $2^{1-2^b}$. Any positive rational separation cutoff
for them has a denominator of at least $2^b$ binary digits. Indeed
$z=\pm x_1=\pm2^{-2^b}$, the weights of~(c) give irreducibility, and a
cutoff $p/q\le2^{1-2^b}$ with $p\ge1$ has $q\ge2^{2^b-1}$, so $q$ has at least $2^b$ binary digits.

% ---------------------------------------------------------------------------
\subsubsection{Zero probes}\label{app:zero-probes}

Without probes, fix a model with exactly one word-against-pair test
$w_0\sim\{w_1,w_2\}$, feasible rational configurations $A$ and~$B$, and
$t_C=\theta_C(w_1)/\theta_C(w_0)$ for $C=A,B$. For a profile~$p$, let
$\Theta_p$ be the set of rational $\theta$ satisfying the probability
floor, normalization and the tests held at~$p$. Let $\Psi$ have one row
$(|w|_m-|w'|_m)_{m\in M}$ for each word-against-word test $w\sim w'$ and
probability-ratio test $\theta(w)=c\,\theta(w')$ held at~$p_A$. Let $f,g$
be the label counts of $w_1,w_2$ minus those of~$w_0$, and put
$(f,g)\colon v\mapsto(f\cdot v,g\cdot v)$. When $(f,g)(\ker_\RR\Psi)$ has
dimension two, let $d_0,d_1,d_\infty$ be the orders of $(1,0),(0,1),(1,1)$
in $\mathbb Z^2/(f,g)(\ker_{\mathbb Z}\Psi)$.

\begin{proposition}[Zero probes: two decidable cases]\label{prop:zero-probes}
In this setting:
\begin{enumerate}
\item[(a)] if $p_A$ does not hold the word-against-pair test, then $B$ is
finely reachable from $A$ if and only if $p_A\preceq p_B$;
\item[(b)] if $p_A$ holds it, and either $(f,g)(\ker_\RR\Psi)$ has
dimension at most one, or it has dimension two and
$1/d_0+1/d_1+1/d_\infty<1$, then $B$ is finely reachable from $A$ if and
only if $p_A\preceq p_B$ and $t_A=t_B$.
\end{enumerate}
On binary inputs, the case distinction and the criterion in~(b) are
polynomial-time.
\end{proposition}

\begin{proposition}[Zero probes: the remaining case]\label{prop:zero-probe-closure}
In the setting of Proposition~\ref{prop:zero-probes}, suppose $p_A$ holds
the word-against-pair test and $(f,g)(\ker_\RR\Psi)$ has dimension two.
Let $\mathcal V=\{\theta(w_1)/\theta(w_0):\theta\in\Theta_{p_A}\}$. Then
$B$ is finely reachable from $A$ if and only if $p_A\preceq p_B$ and
$\overline{\mathcal V}$ contains the closed interval between $t_A$
and~$t_B$.
\end{proposition}

\emph{(Z) Reduction to label probabilities.} %
If $p_A\npreceq p_B$, there is no sequence. Otherwise the target passes the
source tests, so $\theta_B\in\Theta_{p_A}$ and
Theorem~\ref{thm:projected-components}(a) applies. By
Theorem~\ref{thm:projected-components}(a), $B$ is finely reachable from $A$
if and only if $p_A\preceq p_B$ and $A_B$ is finely reachable from~$A$;
these share $p_A$ and~$L_A$. The row sums of the
projected tables give~$\theta$ (Appendix~\ref{app:acquisition-proofs},
(T3)), so a sequence from $A$ to $A_B$ at tolerance $\varepsilon$ gives one
in $\Theta_{p_A}$ from $\theta_A$ to $\theta_B$ with steps below
$N\varepsilon$, $N=|X||R|$. Conversely, by the table formula, a sequence
in $\Theta_{p_A}$ with steps below~$\varepsilon$, kept at $L_A$, is one at
tolerance~$\varepsilon$. So $B$ is finely reachable from $A$ exactly when
$p_A\preceq p_B$ and $\theta_A,\theta_B$ are finely connected
in~$\Theta_{p_A}$, which contains $\theta_B$ as $p_B$ holds the tests
of~$p_A$.

\emph{(G) Homogeneous equations.} Any two points of a set of rational
label-probability vectors cut out by the probability floor, normalization
and equations $\theta(w)=c\,\theta(w')$ with $|w|=|w'|$ and $c>0$ are
finely connected in it, by the proof of Lemma~\ref{lem:transfer}(c).

\begin{proof}[Proof of Proposition~\ref{prop:zero-probes}]
(a) Here $\Theta_{p_A}$ is as in~(G), so (Z) gives the claim.

(b) For $\theta\in\Theta_{p_A}$ put $y_m=\log(\theta_m/\theta_{A,m})$.
Then $\Psi y=0$, and the two ratios $t_Ae^{f\cdot y}$ and $(1-t_A)e^{g\cdot y}$ of $w_1,w_2$
to $w_0$ sum to one. If
$(f,g)(\ker_\RR\Psi)=0$, the ratio is constant. If it is a line spanned by
$(\phi_0,\phi_1)\neq0$, then $(f\cdot y,g\cdot
y)=\varsigma(\phi_0,\phi_1)$ and
$t_Ae^{\phi_0\varsigma}+(1-t_A)e^{\phi_1\varsigma}=1$; the left side is
strictly convex in~$\varsigma$, so the ratio takes at most two values. In
the hyperbolic case, let $P_0$ be the primes dividing a numerator or
denominator of $t_A$ or $1-t_A$, and $\Lambda=(f,g)(\ker_{\mathbb Z}\Psi)$.
The $\ell$-adic valuations of $\theta_m/\theta_{A,m}$ form a vector in
$\ker_{\mathbb Z}\Psi$, so for $\ell\notin P_0$ the valuation pair of
$(t,1-t)$ lies in~$\Lambda$. Write
$t=n_1/n_3$ and $1-t=n_2/n_3$ with $n_1+n_2=n_3$ positive and pairwise
coprime. A prime $\ell\notin P_0$ dividing $n_1$, $n_2$ or $n_3$
contributes $(v_\ell(n_1),0)$, $(0,v_\ell(n_2))$ or
$(-v_\ell(n_3),-v_\ell(n_3))$ to~$\Lambda$, so the exponent is a multiple
of $d_0$, $d_1$ or~$d_\infty$. Hence $n_1=c_1(n'_1)^{d_0}$,
$n_2=c_2(n'_2)^{d_1}$ and $n_3=c_3(n'_3)^{d_\infty}$, where $c_1,c_2,c_3$
are supported on $P_0$ with exponents below the respective orders and
$n'_1,n'_2,n'_3$ share no common factor. Only finitely many coefficient
triples $(c_1,c_2,c_3)$ arise. For each fixed coefficient triple, Darmon
and Granville~\cite[Theorem~2]{darmongranville} give finitely many proper
solutions $(n'_1,n'_2,n'_3)$, since $1/d_0+1/d_1+1/d_\infty<1$.

In both cases the ratio takes finitely many values on~$\Theta_{p_A}$ and is
uniformly continuous, as $\theta(w_0)\ge\delta^{|w_0|}$, so if
$t_A\neq t_B$ small steps cannot change it, and (Z) gives necessity. If
$t_A=t_B$, the vectors of $\Theta_{p_A}$ with $\theta(w_1)=t_A\theta(w_0)$
and $\theta(w_2)=(1-t_A)\theta(w_0)$, which imply the word-against-pair
test, form a set as in~(G) containing $\theta_A$ and~$\theta_B$; (G) and
(Z) give sufficiency.

For the running time, $\dim(f,g)(\ker_\RR\Psi)$ is the rank of $\Psi$
with $f,g$ adjoined minus that of~$\Psi$. Smith forms with multipliers are computed in polynomial
time~\cite[Theorems~4--5 and the rectangular extension following
Theorem~5]{kannanbachem}. %
The right multiplier of the Smith form of $\Psi$ gives an integral kernel
basis. Apply $(f,g)$ to this basis, and take the Smith form of the
resulting two-row generator matrix. Its left multiplier $W$ sends the image
lattice $\Lambda$ to the diagonal lattice with invariant factors
$s_1,s_2$, so $\mathbb Z^2/\Lambda\cong\mathbb Z/s_1\times\mathbb Z/s_2$,
and the order of $e$ is
$\operatorname{lcm}_i\bigl(s_i/\gcd(s_i,(We)_i)\bigr)$. The ratios
$t_A,t_B$ are compared exactly, and no step factors a prime.
\end{proof}

\begin{proof}[Proof of Proposition~\ref{prop:zero-probe-closure}]
The probability floor gives $\theta(w_1),\theta(w_2)\ge\delta^{|w_0|}$,
and $\theta(w_0)\le1$, so both $t$ and $1-t$ are uniformly bounded away
from zero. By~(Z) it suffices to show that $\theta_A$ and $\theta_B$ are
finely connected in $\Theta_{p_A}$ exactly when $\overline{\mathcal V}$
contains the interval $I$ between $t_A$ and~$t_B$. Fix a label $m_0$, let
$\mathcal Z=\{z\in\RR^M:z_{m_0}=0,\ \Psi z=0\}$, and map $z$ to the
normalization of $(\theta_{A,m}e^{z_m})_m$, a homeomorphism onto the
normalized positive real vectors satisfying the held homogeneous tests. The
probability floor becomes the convex set
\[
  K=\Bigl\{z\in\mathcal Z:\ \sum_\ell\frac{\theta_{A,\ell}}{\theta_{A,m}}\,e^{z_\ell-z_m}\le\frac1\delta\ \text{ for every }m\Bigr\},
\]
which is compact. On the open segment between distinct points of $K$,
every inequality is strict: for each~$m$, some $z_\ell-z_m$ differs
between the endpoints, which agree at~$m_0$, and that term is strictly
convex. Put $\pi(z)=(f\cdot z,g\cdot z)$, which maps $\mathcal Z$ onto
$\RR^2$: every defining vector of $\Psi$, as well as $f,g$, annihilates
constants, so fixing $z_{m_0}=0$ leaves the rank-two image unchanged. Also
put $\gamma(t)=(\log(t/t_A),\log((1-t)/(1-t_A)))$. The
word-against-pair test says $\pi(z)=\gamma(t)$ with $t=t_Ae^{f\cdot z}$. The
closure $F$ in $K$ of the image of $\Theta_{p_A}$ is compact, and the
ratio maps it onto $\overline{\mathcal V}$.

\emph{Full fibers.} For $t\in\overline{\mathcal V}$, the fiber
$F_t=F\cap\pi^{-1}(\gamma(t))$ is the whole real fiber
$K_t=K\cap\pi^{-1}(\gamma(t))$. If $\gamma(t)$ lies on the boundary of
$\pi(K)$, then $K_t$ has at most one point, since the midpoint of two would
be interior to $K$ and $\pi$ is open; and $F_t\neq\emptyset$.
If $\gamma(t)$ is interior to $\pi(K)$, then $K$ has interior, since a
midpoint of two distinct points satisfies every probability-floor
inequality strictly. Start with an interior point $z_1$. If
$\pi(z_1)=\gamma(t)$, it is the desired lift. Otherwise extend the
projected segment slightly beyond $\gamma(t)$ within $\pi(K)$, lift its far
endpoint and interpolate back. Convex interpolation with this interior lift
makes interior points dense in~$K_t$. For such an interior $z^*\in K_t$
and $t_j\in\mathcal V$ with $t_j\to t$, a fixed linear right inverse $\rho$
of $\pi$ gives, for large~$j$, interior points
$z^*+\rho(\gamma(t_j)-\gamma(t))$ over $\gamma(t_j)$. Take
$\theta^{(j)}\in\Theta_{p_A}$ with ratio~$t_j$. The real fiber of
$\mathcal Z$ over $\gamma(t_j)$ is the log image of $\theta^{(j)}$ plus a
rational subspace, the common kernel of $\Psi$, $f$, $g$ and $z_{m_0}$, with
integral basis $u_1,\dots,u_r$. Multiplying $\theta^{(j)}$ by
$\prod_iq_i^{u_i}$ with positive rationals $q_i$ and normalizing keeps every held equation and the ratio. In log
coordinates the twists add $\sum_i\log(q_i)u_i$. %
Since $\log\mathbb Q_{>0}$ is dense in~$\RR$, for each sufficiently large~$j$
choose a twist whose log image is inside $K$ and within $1/j$ of the
translated interior point. These log images come from $\Theta_{p_A}$ and
converge to~$z^*$. Thus $z^*\in F$, and closedness of $F$ gives
$K_t\subseteq F$ by density of the interior lifts in~$K_t$.

\emph{Components.} Each $F_t$ is convex, hence connected. If
$I\subseteq\overline{\mathcal V}$, the compact set of points of $F$ with
ratio in $I$ is connected, since each fiber is connected and a separation would give two disjoint compact ratio images covering~$I$. Conversely,
the ratio maps a connected subset of $F$ containing the images of
$\theta_A$ and $\theta_B$ onto an interval containing~$I$. So the two
images share a component of $F$ exactly when
$I\subseteq\overline{\mathcal V}$, and Lemma~\ref{lem:compact-chain},
applied to $\Theta_{p_A}$ in its compact closure, finishes the proof.
\end{proof}

% ---------------------------------------------------------------------------
\subsubsection{The one-state subclass}\label{app:one-state}

% Sources: ComplexityBoundary.md section 2.2 (cubical realization,
% compilation, tractable restrictions) and section 3 (gate normalization,
% gated box, templates). Scoped symbols: the cube set Q_cube, counts
% n_v, n_c, n_e, k_0, the node bound Omega, K (delta = 1/K).
\begin{proof}[Proof of Proposition~\ref{prop:one-state}]
(a) \emph{Source.} Deciding whether two satisfying assignments of a 3-CNF
formula are joined by single-variable flips through satisfying assignments
is PSPACE-complete~\cite[Lemma~3.6]{gopalan}. Map instances with equal
assignments to a fixed yes-instance with two distinct adjacent satisfying
assignments, and delete tautological clauses. Let $\mathcal Q$ be the set
of rational points of $[0,1]^{n_{\mathrm v}}$ at which every clause has a
true literal, $x_i$ being true at $x_i=1$ and $\neg x_i$ at $x_i=0$. For
$x\in\mathcal Q$, every vertex of the smallest closed face containing $x$
satisfies the formula.
Conversely, if a clause has no true fixed literal on a face, its remaining
literals can all be made false at a corner, because the clause is
non-tautological. Thus $\mathcal Q$ is exactly the set of rational points
in the union of faces of Schnider and Weber's induced cubical
complex~\cite[Definition~4]{schniderweber}. The rational points on the segment of a satisfying flip lie in~$\mathcal Q$. Two points of $\mathcal Q$ at sup distance
below one fix no coordinate to opposite values, so their smallest faces
share a vertex; all corners of a face are flip-connected. So steps below
one preserve the Boolean flip class, and flip-connected assignments are
joined at every tolerance.

\emph{Compiler.} Let $\delta=1/K$.
Use maximum word length $H=2$ and give each emitted test its own name.
The sole admissible word-test set is the full catalogue. Use no value tests
and give both endpoints the full-support, level-zero profile holding the
whole catalogue. With singleton base and register and positive label
probabilities, every feasible configuration has that support and register
partition. Thus the compiled label-probability vectors are exactly the
feasible configurations at the common profile.
The labels are two references of probability $2^j\delta$ at each level
$j=0,\dots,k_0$, a pair of probabilities $\delta(1+2x_i)$, $\delta(3-2x_i)$
per coordinate, and, per distinct node of the clause expressions, two
copies of an output--complement pair with sum $2^{k_0}\delta$. %
Pad all rate equalities and sums with the first level-zero reference~$r$.
Equate the two references at each level, and for each positive level
equate its first reference to the sum of the two references at the
preceding level. Equate each coordinate-pair sum to a level-two reference
and each output--complement pair sum to a level-$k_0$ reference. For every
expression node, equate its two output labels and impose its defining
equation on the first output label. A constant uses~$r$, and a leaf uses
its coordinate label. A sum uses the first output copy of the left child
and the second of the right child, so its two right-hand words are
distinct even for repeated children. For a product with output~$z$ and
child outputs $a,b$, compare $zr$ with~$ab$.
For each clause, equate the output labels of its two positive expression
trees: the product of the factors $(3-2x_i)-1$ for its positive literals
and $(1+2x_i)-1$ for its negative ones, written as a difference of positive trees by $(P-Q)(P'-Q')=(PP'+QQ')-(PQ'+QP')$. Their difference is the product of the indicated literal factors. %
The bound of a node is its value when every coordinate leaf is replaced
by~$3$ and constant leaves are kept at~$1$. Start a clause product with the
pair $(1+1,1)$. Each additional literal multiplies the maximum bound by at
most~$4$ and changes the total size $s$ of the two trees to $2s+10$.
Starting at bound $2$ and size $4$, width at most three gives bound $128$
and at most $102$ nodes. Take $\Omega$ to be the maximum of zero and all
node bounds, and put $k_0=\max(2,\lceil\log_2(\Omega+1)\rceil)$; for width
three, $\Omega\le128$, $k_0\le8$, and $n_{\mathrm c}$ clauses have
$n_{\mathrm e}\le102n_{\mathrm c}$ nodes.
There are $2(k_0+1)+2n_{\mathrm v}+4n_{\mathrm e}$ labels,
$2k_0+1+n_{\mathrm v}+4n_{\mathrm e}+n_{\mathrm c}$ word tests and
$K=2(2^{k_0+1}-1)+4n_{\mathrm v}+2^{k_0+1}n_{\mathrm e}$. The reference
equations and normalization make the level-$0$ probability~$\delta$, and
induction over the nodes, whose values lie in $[1,2^{k_0}-1]$, gives a
bijection between $\mathcal Q$ and the feasible configurations. As
$x_i=(\theta/\delta-1)/2$ for the first label of its pair, a step below
$2\delta\tau$ moves the cube point by less than~$\tau$, and the inverse is
Lipschitz; so tolerances at most $2\delta$ preserve the flip class. %
The lift is polynomial in the cube coordinates and hence Lipschitz on the
compact cube, so rational subdivisions of satisfying flip segments lift at
every positive tolerance. At Boolean endpoints every label probability is
an integer multiple of $1/K$. Since $k_0$ is bounded,
$K=O(1+n_{\mathrm v}+n_{\mathrm c})$. Because each test uses a constant
number of label names and one entry lists the full catalogue as the
admissible set, the output has binary length
$O((1+n_{\mathrm v}+n_{\mathrm c})\log(2+n_{\mathrm v}+n_{\mathrm c}))$
and is written in polynomial time.

(b) Length-one tests give linear equations. Linear held equations make the
feasible set convex, so rational points of the segment from $\theta_A$ to
$\theta_B$ give sequences at every tolerance. If every held word test
is word-against-word, (G) of Appendix~\ref{app:zero-probes} applies. With
one base state and one register value, the nonconvex example in
Section~\ref{sec:zero-probes} has exactly two feasible label-probability
vectors, so every tolerance at most their distance blocks a switch.

(c) In cube coordinates, the switch construction of
Section~\ref{sec:rational-roots} has $s=2x_s$, $\zeta=2x_\zeta$,
$z_\gamma=2x_\gamma-1$, $\psi=2x_\psi-1$ and $\xi=1+2x_\xi$. Each
gated residual is a difference of positive trees in the leaves $1$,
$A_i=1+2x_i$ and $3-2x_i$, compiled as in~(a).
Use $(1+1,1)$ for $1$, $(A_i,1)$ for $2x_i$, coordinatewise addition, and
pair swap for negation. Associate the slack residual as
$\xi^2-(2+(\zeta\zeta)\psi)$. A pair product adds at most two to the
maximum depth, and a sum adds at most one. Pair bounds add under sums and
at most double the product under products. The pairs for
$z_\gamma,\psi,\xi$ have bound at most~$5$, those for $s,\zeta$ at
most~$3$, and that for $2$ at most~$4$. The switch factor has bound at
most $2\cdot3(4+3)=42$, each residual at most $234$, and each gated output
at most $2\cdot42\cdot234=19656$. Trees have depth at most $8$ and at most
$1022$ nodes, so $k_0\le15$ and $K=O(\mathcal L+1)$. The endpoints are the vertices
with $x_s\in\{0,1\}$ and all other coordinates~$0$, feasible since the
switch factor vanishes. If $f$ has a rational zero, with gate
values~$v_\gamma$, put $\zeta=(1+\sum_\gamma|v_\gamma|)^{-1}$,
$z_\gamma=\zeta v_\gamma$, $\xi$ rational in
$(\sqrt2,\sqrt2+\zeta^2/4)$ and $\psi=(\xi^2-2)/\zeta^2$, so $1<\xi<2$ and
$0<\psi<1$. Every residual then vanishes for every~$s$ and the end faces
are full boxes, so lifting three rational segments joins the endpoints at
every tolerance. Otherwise no feasible
rational point has $0<s<2$ (Section~\ref{sec:rational-roots}), and a step
from $x_s=0$ to $x_s=1$ changes a table by at least $2\delta$. Over $\RR$,
the point $\zeta=0$, $z_\gamma=0$, $\psi=0$, $\xi=\sqrt2$ satisfies the
gated equations for every~$s$. A real feasible point with $0<s<2$ reaches an end face by varying $s$ alone. The end faces are boxes, so the real gated set is path-connected. The compiler's induction works over the reals as well, so its continuous image, the real feasible set, is path-connected. The output is polynomial, as in~(a), with $O(\mathcal L+1)$ gates and a fixed-size template per residual.
\end{proof}
